% !TeX root = ../../lagrangian-torus-fibration.tex
% 2026-11-20 / v1.0: 全閉 triple corner の GH 極限と fiber-preserving relative averaging．
\section{全閉 triple corner の内在的 GH 極限}
\label{ke:closed-corner-gh-section}

系\ref{ke:actual-closed-corner}で構成した actual Lagrangian bundle に対し，
本節では三枚の閉 square の共通 face を含む領域全体の内在距離を同定する．
Lagrangian 化に用いた Moser flow と fibers は同じものを使う．
新たな距離の入力は，関連 KE 主稿\cite[Part~IV,
\texttt{simplex:thm:closed-corner-model}]{KELocalCurrent}の
complete simplex model に関する全閉 corner の評価である．
同稿の \texttt{simplex:cor:actual-closed-corner-metric} は，
Moser 変形前の actual 領域について対応する結論を与える．
以下では model metric 自体の Lie derivative を評価し，
既に構成した Moser 像へ sharp な距離比較を移す．
さらに，三重点の外側保護領域内を通る経路を全て許した距離についても，
同じ projection による極限を示す．

\subsection{Closed corner 自身の距離と model の入力}

前節と同じ任意の固定 $d\ge5$ の族を取り，その genericity 仮定を保つ．
三重点の label を $J=\{i,j,k\}$，
$L=-\log|t|^2$，$R=\log L$，$m=R$ とする．
\[
 \mathcal C_J
 =\{w\in[0,1]^J:\max_{a\in J}w_a=1\}
\]
に，三枚の unit Euclidean square を共有 faces に沿って貼り合わせた
\emph{この corner 自身の内在長さ距離} $d_J$ を入れる．
この記号は，$\mathcal C_d$ 全体の距離 $d_\square$ を
$\mathcal C_J$ に制限したものを意味しない．

命題\ref{ke:triple-hexagon}の softmax を
\[
 \ell_a(w)=R+(L-3R)\frac{L^{w_a}}{\sum_{b\in J}L^{w_b}}
\]
とし，その六角形 $H_{L,R}$ の全 phase torus preimage を
\[
 \mathcal H_t=\mathcal U_t(H_{L,R}),\qquad
 F_{t,R}:\mathcal C_J\times\TT^2\longrightarrow\mathcal H_t
\]
と書く．$F_{t,R}$ は既存の parameter 化と同じであり，
\[
 P_t=\operatorname{pr}_{\mathcal C_J}\circ F_{t,R}^{-1}
\]
は変形前の softmax projection である．
Complete simplex KE form $k_L$ の実計量を $\gamma_t=g_L^{\mathrm{sf}}$ と書き，
$\mathcal H_t$ の中に留まる曲線の $R^{-2}\gamma_t$-長さで定義した
内在距離を $d_t^{\mathrm{sf,int}}$ とする．

上記 KE 主稿の \texttt{simplex:thm:closed-corner-model} は
\begin{equation}\label{ke:eq:closed-corner-model-input}
 \sup_{z,z'\in\mathcal H_t}
 \left|d_t^{\mathrm{sf,int}}(z,z')
       -d_J(P_t(z),P_t(z'))\right|
 \le C\frac{\log R}{R}
\end{equation}
を与える．この評価は全 closed corner と全 phase に一様である．
特に $P_t$ の全射性から
\begin{equation}\label{ke:eq:closed-model-diameter}
 \operatorname{diam}(\mathcal H_t,d_t^{\mathrm{sf,int}})
 \le\operatorname{diam}(\mathcal C_J,d_J)+C\frac{\log R}{R}
\end{equation}
は一様有界である．以下の actual estimates の定数には
固定した $d$，族，Hermitian data，chart への依存を許す．
次数や族全体にわたる一様性は主張しない．

\subsection{同じ Moser flow による sharp な metric 比較}

\begin{lemma}[全閉 corner の Moser pullback metric]
\label{ke:closed-corner-moser-metric}
系\ref{ke:actual-closed-corner}の flow $\phi_{t,s}$ と
$\phi_t=\phi_{t,1}$ を用いる．十分小さい $t\ne0$ に対し，
\begin{equation}\label{ke:eq:closed-model-flow-metric}
 e^{-C/R}\gamma_t\le\phi_{t,s}^*\gamma_t
                 \le e^{C/R}\gamma_t
 \qquad(0\le s\le1)
\end{equation}
が $\mathcal U_t(D_{L,R})$ 上で成立する．さらに
\begin{equation}\label{ke:eq:closed-actual-pullback}
 (1-C/R)\gamma_t\le\phi_t^*g_t\le(1+C/R)\gamma_t
\end{equation}
となる．
\end{lemma}
\begin{proof}
$h=h_t$ を\eqref{ke:eq:cusp-model}の constrained cusp metric とする．
関連 KE 主稿の complete simplex model の一様幾何と
\texttt{coarse:thm:full-triple-tensor} により，外側領域
\[
 \mathcal V_t=\mathcal U_t(\ell_a>R/2\text{ for all }a\in J)
\]
上で
\begin{equation}\label{ke:eq:closed-model-c1-data}
 c h\le\gamma_t\le C h,
 \qquad |\nabla^h\gamma_t|_h\le C,
 \qquad |g_t-\gamma_t|_h\le C/R
\end{equation}
が成り立つ．深度 $m=R/2$ で同定理を用いている．
また\eqref{ke:eq:full-corner-moser-c1}より
\[
 |Y_{t,s}|_h+|\nabla^hY_{t,s}|_h\le C/R
\]
であり，内側領域からの flow は $\mathcal V_t$ に留まる．

任意の接ベクトル $U,V$ に対し，$\nabla^h$ の torsion が零であることから
\[
 (\mathcal L_{Y_{t,s}}\gamma_t)(U,V)
 =(\nabla^h_{Y_{t,s}}\gamma_t)(U,V)
  +\gamma_t(\nabla^h_UY_{t,s},V)
  +\gamma_t(U,\nabla^h_VY_{t,s}).
\]
従って\eqref{ke:eq:closed-model-c1-data}より
\[
 |\mathcal L_{Y_{t,s}}\gamma_t|_h\le C/R,
 \qquad
 -\frac CR\gamma_t\le\mathcal L_{Y_{t,s}}\gamma_t
                         \le\frac CR\gamma_t.
\]
$\frac{d}{ds}\phi_{t,s}^*\gamma_t
=\phi_{t,s}^*(\mathcal L_{Y_{t,s}}\gamma_t)$ に
Gronwall の不等式を用いると\eqref{ke:eq:closed-model-flow-metric}を得る．
さらに $\gamma_t\ge ch$ なので，最後の metric difference は
\[
 (1-C/R)\gamma_t\le g_t\le(1+C/R)\gamma_t
 \quad\text{on }\mathcal V_t
\]
を意味する．これを $\phi_t$ で引き戻し，
\eqref{ke:eq:closed-model-flow-metric}と合わせると，定数を増やして
\eqref{ke:eq:closed-actual-pullback}を得る．
背景間の coarse comparison だけで sharp comparison を置き換えず，
model $\gamma_t$ の Lie derivative を直接用いた点がここで必要である．
\end{proof}

\begin{theorem}[Actual Lagrangian bundle による全閉 corner の内在的 GH 収束]
\label{ke:closed-corner-gh}
前節の系\ref{ke:actual-closed-corner}と同じ
\[
 \mathcal W_{t,J}=\phi_t(\mathcal H_t),\qquad
 \pi_{t,J}=P_t\circ\phi_t^{-1}
\]
を取る．$\mathcal W_{t,J}$ の中に留まる piecewise smooth curves の
$R^{-2}g_t$-長さから定まる内在距離を $d_{t,J}^{\mathrm{Lag,int}}$ とする．
すると
\begin{equation}\label{ke:eq:closed-corner-distortion}
 \sup_{x,y\in\mathcal W_{t,J}}
 \left|d_{t,J}^{\mathrm{Lag,int}}(x,y)
      -d_J(\pi_{t,J}(x),\pi_{t,J}(y))\right|
 \le C\frac{\log R}{R}.
\end{equation}
従って同じ actual $\omega_t$-Lagrangian $T^2$ fibration の射影
$\pi_{t,J}$ 自身が全射 GH approximation となり，
\begin{equation}\label{ke:eq:closed-corner-gh-limit}
 (\mathcal W_{t,J},d_{t,J}^{\mathrm{Lag,int}})
 \xrightarrow[t\to0]{\mathrm{GH}}(\mathcal C_J,d_J)
\end{equation}
を得る．
全 fiber の非正規化の intrinsic diameter は $O(R^{-1})$，
正規化後は $O(R^{-2})$ であり，
raw map との\eqref{ke:eq:actual-closed-corner-target}も保たれる．
\end{theorem}
\begin{proof}
補題\ref{ke:closed-corner-moser-metric}を $\mathcal H_t$ に制限する．
$\phi_t$ はその開近傍で像への diffeomorphism であるから，
$\mathcal H_t$ 内の経路と $\mathcal W_{t,J}$ 内の経路は一対一に対応する．
従って $\varepsilon_t=C/R<1/2$ と置けば
\[
 (1-\varepsilon_t)^{1/2}d_t^{\mathrm{sf,int}}(z,z')
 \le d_{t,J}^{\mathrm{Lag,int}}(\phi_t(z),\phi_t(z'))
 \le(1+\varepsilon_t)^{1/2}d_t^{\mathrm{sf,int}}(z,z').
\]
\eqref{ke:eq:closed-model-diameter}により model diameter は一様有界なので，
この差は一様に $O(R^{-1})$ である．
\eqref{ke:eq:closed-corner-model-input}と
$\pi_{t,J}\circ\phi_t=P_t$ を合わせると
\eqref{ke:eq:closed-corner-distortion}を得る．

$P_t$，従って $\pi_{t,J}$ は全射である．その graph を correspondence として用いれば，
\eqref{ke:eq:closed-corner-distortion}から GH 距離の上界と
\eqref{ke:eq:closed-corner-gh-limit}が従う．
Fiber と raw map に関する記述は，bundle と Moser flow を変更していないので
系\ref{ke:actual-closed-corner}そのものである．
\end{proof}

\begin{corollary}[外側保護領域を通る距離に関する同じ極限]
\label{ke:closed-corner-protected-gh}
定理\ref{ke:closed-corner-gh}と同じ $\mathcal W_{t,J}$ と $\pi_{t,J}$ を用いる．
外側保護領域
\[
 \mathcal V_t
 =\{\zeta_i\zeta_j\zeta_k=t,\ \min_{a\in J}\ell_a>R/2\}
\]
内の全 piecewise smooth curves の $R^{-2}g_t$-長さの下限で定まる距離を，
$\mathcal W_{t,J}$ へ制限して $d_{t,J}^{\mathrm{Lag,out}}$ と書く．
すると
\begin{equation}\label{ke:eq:closed-corner-protected-distortion}
 \sup_{x,y\in\mathcal W_{t,J}}
 \left|d_{t,J}^{\mathrm{Lag,out}}(x,y)
      -d_J(\pi_{t,J}(x),\pi_{t,J}(y))\right|
 \le C\frac{\log R}{R}.
\end{equation}
特に $\pi_{t,J}$ 自身によって
\[
 (\mathcal W_{t,J},d_{t,J}^{\mathrm{Lag,out}})
 \xrightarrow[t\to0]{\mathrm{GH}}(\mathcal C_J,d_J)
\]
となる．経路は $\mathcal W_{t,J}$ の外へ出てもよいが，
$\mathcal V_t$ の外へは出ないものに限る．
\end{corollary}
\begin{proof}
外側保護領域には全ての Moser trajectories が入るので，
$\mathcal W_{t,J}\subset\mathcal V_t$ である．従って
\[
 d_{t,J}^{\mathrm{Lag,out}}\le d_{t,J}^{\mathrm{Lag,int}}
\]
となり，上界は定理\ref{ke:closed-corner-gh}から従う．

下界には，関連 KE 主稿\cite[Part~IV,
\texttt{simplex:lem:closed-corner-no-shortcut}]{KELocalCurrent}の
両端を clipped した projection を用いる．
$[s]_0^1=\max\{0,\min\{s,1\}\}$ と置き，
\[
 S=R-2\log R,\qquad
 \Psi_a=\left[\frac{\log(\ell_a/R)}S\right]_0^1
 \quad(a\in J)
\]
と定める．これは complete simplex 全体から $\mathcal C_J$ への写像である．
同補題は，complete model 内の任意の曲線 $c$ に対して
\begin{equation}\label{ke:eq:protected-pruning-length}
 \operatorname{Length}_{d_J}(\Psi\circ c)
 \le\frac{\sqrt{1+CR^{-2}}}{S}\operatorname{Length}_{\gamma_t}(c)
\end{equation}
を与え，また
\begin{equation}\label{ke:eq:protected-pruning-endpoints}
 \sup_{z\in\mathcal H_t}d_J(\Psi(z),P_t(z))
 \le C\frac{\log R}{R}
\end{equation}
を与える．この長さ評価には clipping level と square seams 上の経路も含まれる．

特に $z\in\mathcal H_t$ の Moser trajectory に対し，
$|Y_{t,s}|_{\gamma_t}\le C/R$ と
\eqref{ke:eq:protected-pruning-length}から
\begin{equation}\label{ke:eq:protected-moser-endpoint}
 d_J(\Psi(\phi_t(z)),\Psi(z))
 \le\frac{C}{S}\int_0^1|Y_{t,s}(\phi_{t,s}(z))|_{\gamma_t}\,ds
 \le\frac{C}{R^2}
\end{equation}
を得る．従って $x=\phi_t(z)$ について
\[
 d_J(\Psi(x),\pi_{t,J}(x))\le C\frac{\log R}{R}.
\]

さらに\eqref{ke:eq:closed-model-c1-data}による
$(1-C/R)\gamma_t\le g_t$ は $\mathcal V_t$ 全体で成立する．
$\mathcal V_t$ 内の $x,y$ を結ぶ任意の曲線に
\eqref{ke:eq:protected-pruning-length}を用いて下限を取ると
\[
 d_{t,J}^{\mathrm{Lag,out}}(x,y)
 \ge \frac{S\sqrt{1-C/R}}{R\sqrt{1+CR^{-2}}}
          d_J(\Psi(x),\Psi(y)).
\]
右辺の係数は $1-O(\log R/R)$ であり，$d_J$ の diameter は有限である．
端点評価を加えると
\[
 d_{t,J}^{\mathrm{Lag,out}}(x,y)
 \ge d_J(\pi_{t,J}(x),\pi_{t,J}(y))-C\frac{\log R}{R}
\]
となる．既に示した上界と全射性から，distortion bound と GH 極限を得る．
\end{proof}

\begin{remark}[Corner 内の二つの距離と大域距離]
\label{ke:closed-corner-gh-scope}
定理\ref{ke:closed-corner-gh}は三枚の square を合わせた corner 全体を扱うが，
曲面全体の距離を扱うものではない．$d_{t,J}^{\mathrm{Lag,int}}$ には
$\mathcal W_{t,J}$ 内の経路を用い，
系\ref{ke:closed-corner-protected-gh}の $d_{t,J}^{\mathrm{Lag,out}}$ には
$\mathcal V_t$ 内の経路を全て許す．
$X_t$ 内で $\mathcal V_t$ の外へ出る経路は，いずれにも含めていない．
また標的の $d_J$ も $\mathcal C_J$ 内の経路だけを用いる．
Full complex $\mathcal C_d$ の他の cells を通る経路との比較は，
本節のいずれの distortion bound にも含めていない．

従って異なる三重点 chart の接合，neck・source・branch・core を含む
既存の大域 fibration への相対接着，および大域 shortcut の評価は残る．
本節の regular fibers と既存大域 fibration の singular fibers を同一視しない．
解析入力の新規性の包括的確認も引き続き必要である．
\end{remark}

\subsection{次の相対接着への条件}

Fibration の接着に必要な collar の固定と，各点を固定することを区別する．
以下の命題は，既存の averaging--Moser 構成に直接適用できる．
Actual neck にその仮定を満たす collar が存在することを主張するものではない．

\begin{proposition}[既に Lagrangian な orbit collar を保つ averaging]
\label{ke:relative-fiber-preserving}
定理\ref{ke:short-torus-averaging}の設定で，
$N\subset\mathcal U'$ を torus 飽和した開集合とする．
$N$ 内の元の全 phase tori が既に $\omega$-Lagrangian であると仮定する．
同定理で用いた chain homotopy primitive
\[
 \beta=H\omega,
 \qquad H=K_1+\mathcal A_1K_2
\]
と $\omega_s=(1-s)\bar\omega+s\omega$ による Moser flow は，
$N$ の各 phase torus を集合として保つ．
従って $\phi_s(N)=N$ であり，完成した fibration は
\begin{equation}\label{ke:eq:relative-fibration-identity}
 (p\circ\phi_1^{-1})|_N=p|_N
\end{equation}
を満たす．$\omega$ の torus 不変性は仮定しない．
同じ結論は remark~\ref{ke:relative-averaging}の centered chain homotopy にも成立する．
\end{proposition}
\begin{proof}
Vertical tangent distribution を $\mathcal E=\ker dp$ と書く．
仮定より $\omega|_{\mathcal E\times\mathcal E}=0$ on $N$ であり，
torus の各 translation は $N$ と $\mathcal E$ を保つ．
従って，生成場 $V_j$ と任意の vertical vector $Z$ について
\[
 \bigl(\rho_{j,s}^*(\iota_{V_j}\omega)\bigr)(Z)
 =\omega_{\rho_{j,s}(x)}
     \bigl(V_j,\,d\rho_{j,s}(Z)\bigr)=0.
\]
積分と他の円作用による平均を取ってもこの消滅は保たれるので
\begin{equation}\label{ke:eq:relative-primitive-vertical}
 \beta|_{\mathcal E}=0\quad\text{on }N.
\end{equation}
Centered homotopy でも積分の weight が変わるだけで，同じ議論が成立する．

$\bar\omega$ も $\mathcal E$ 上で零であり，
$\omega_s$ は K\"ahler form の凸結合で非退化である．
従って実次元二の $\mathcal E$ は各 $\omega_s$ に関して Lagrangian である．
$\iota_{Y_s}\omega_s=\beta$ と
\eqref{ke:eq:relative-primitive-vertical}より
\[
 \omega_s(Y_s,Z)=0\quad(Z\in\mathcal E),
 \qquad Y_s\in\mathcal E^{\omega_s}=\mathcal E
 \quad\text{on }N.
\]
したがって各 trajectory の基底値は一定である．
各 orbit は compact torus なので，制限された時間依存 vector field の flow は
正負両方向に全区間で存在し，各 orbit をそれ自身へ微分同相に写す．
全体の flow との一意性から $\phi_s(N)=N$ および
\eqref{ke:eq:relative-fibration-identity}が従う．
\end{proof}

\begin{corollary}[同じ orbit partition を持つ neck collar との局所接着]
\label{ke:relative-neck-matching}
上の $N$ が別の領域の actual $\omega$-Lagrangian fibration と重なり，
その fibration の fibers が $N$ 上で $p$ の phase tori と一致すると仮定する．
両基底の collar coordinates をこの orbit partition で同定すれば，
averaging 後の fibration と与えられた fibration は $N$ 上で一致し，
通常の開集合上の写像の接着によって一つの局所 fibration となる．
\end{corollary}
\begin{proof}
与えられた fibers の Lagrangian 性により
命題\ref{ke:relative-fiber-preserving}を適用できる．
同命題は基底値も含めて\eqref{ke:eq:relative-fibration-identity}を与えるので，
指定された基底 collar の同定後には二つの写像が open overlap 上で等しい．
その写像と局所 trivializations を接着すればよい．
ここでの結論は重なり近傍の局所的なものであり，
基底全体の Hausdorff 性や大域 properness を別の仮定なしに主張しない．
\end{proof}

\begin{remark}[点ごとの固定と toric 座標変更の区別]
\label{ke:relative-neck-obstructions}
命題\ref{ke:relative-fiber-preserving}は fiber を集合として固定し，
各点を固定するとは述べない．$\phi_s|_N=\id$ と
$\phi_s^*\omega_s=\bar\omega$ を同時に要求するなら，
$\omega_s|_N=\bar\omega|_N$ が必要である．
その意味で，不変 collar を仮定する remark~\ref{ke:relative-averaging}の
点ごとの固定とは異なる結論である．

また，独立な二つの toric coordinates の変更が
$\zeta'=c(t)\zeta^A$，$A\in\mathrm{GL}(2,\ZZ)$ という定数係数の monomial 変更なら，
phase tori の orbit partition は変わらない．
しかし非定数の holomorphic unit を含む
$\zeta'=\zeta^A u(\zeta,t)$ という変更は，一般にその partition 自体を変える．
曲面の同じ overlap に二つの torus actions が存在することだけでは，
系\ref{ke:relative-neck-matching}の仮定は満たされない．
本稿の actual neck と三重点領域の間で，必要な orbit partition，
action periods，および quantitative buffer を同時に合わせることは残る課題である．
上の fiber-preserving 証明では chain homotopy primitive を用いた．
同じ exterior derivative を持つ任意の別の primitive へ替えても，
\eqref{ke:eq:relative-primitive-vertical}が自動的に保たれるわけではない．
\end{remark}
